% !PRINTED 495
\weberpage{495}

%<<CONT
betrachtet. Einen davon verschiedenen Weg hat
\gesperrt{Kronecker}\textsuperscript{1)} eingeschlagen.

Die Theorie von \gesperrt{Dedekind} ist von ihrem Begr\"under umfassend und
in stets wachsender Einfachheit und Vollkommenheit dargestellt in dem
letzten Supplement der drei neuesten Auflagen von \gesperrt{Dirichlet's}
Vorlesungen \"uber Zahlentheorie.

Die Theorie von \gesperrt{Kronecker} ist erst im Jahre 1882 durch die
Festschrift zu \gesperrt{Kummer's} Jubil\"aum dem weiteren Kreise der
Mathematiker bekannt geworden, und ist auch jetzt noch nicht ganz leicht
zug\"anglich\textsuperscript{2)}.

Bez\"uglich der \gesperrt{Dedekind'schen} Theorie k\"onnen wir also den
Leser, der tiefer in diesen interessanten Gegenstand einzudringen
w\"unscht, auf die genannten Werke verweisen. Dagegen soll hier der Versuch
gemacht werden, den Zusammenhang zwischen den verschiedenen Theorieen
herzustellen und auf Grund der \gesperrt{Kronecker'schen} Voraussetzungen
und H\"ulfsmittel so viel daraus zu entwickeln, als f\"ur die folgenden
algebraischen Anwendungen nothwendig ist.

\begin{center}
\S.\,135.

\gesperrt{Ganze Functionen in einem algebraischen K\"orper.}
\end{center}

Schon im ersten Bande haben wir mehrfach Gelegenheit gehabt, ganze
Functionen einer beliebigen Anzahl von unabh\"angigen Ver\"anderlichen
einzuf\"uhren, und haben auch (in \S.\,141 f.)
%CONT>>

\vspace{1em}\hrule\vspace{0.5em}

\noindent 1878). Ueber die Discriminanten endlicher K\"orper (ebend.,
Bd.\,29, 1882). Ueber die Anzahl der Ideal-Classen in den verschiedenen
Ordnungen eines endlichen K\"orpers (Braunschweig 1877). Festschrift zur
S\"acularfeier des Geburtstages von \gesperrt{Gauss}. \glqq Zur Theorie der
Ideale\grqq{} und \glqq Ueber die Begr\"undung der Idealtheorie\grqq{}.
Nachrichten d. Ges. d. Wissensch. in G\"ottingen 1894, 1895. Hierher
geh\"oren auch die Abhandlungen von \gesperrt{Hilbert}, \glqq Ueber die
Zerlegung der Ideale etc.\grqq{}, Mathem. Annalen, Bd.\,44, 1893. \glqq
Grundz\"uge einer Theorie der \gesperrt{Galois}'schen Zahlk\"orper\grqq{}, G\"ottinger
Nachrichten 1894. \gesperrt{Hurwitz}, \glqq Zur Theorie der Ideale\grqq{}.
\glqq Ueber einen Fundamentalsatz etc.\grqq{}, G\"ottinger Nachrichten 1894,
1895.

\noindent\textsuperscript{1)} \gesperrt{Kronecker}, Grundz\"uge einer
arithmetischen Theorie der algebraischen Gr\"ossen. Festschrift zu
\gesperrt{Kummer's} 50j\"ahrigem Doctor-Jubil\"aum. Berlin 1882. (Auch in
Bd.\,92 von \gesperrt{Crelle}'s Journal.) Zu erw\"ahnen sind hier noch die Arbeiten von
\gesperrt{Hensel} in den B\"anden 101, 103, 105, 111, 113 des \gesperrt{Crelle}'schen
Journals.

\noindent\textsuperscript{2)} Vgl. die Vorrede zur 4. Auflage von
\gesperrt{Dirichlet-Dedekind's} Vorlesungen \"uber Zahlentheorie.
